\TieudeT{PHONG TỎA VẬT LÍ 11 - DAO ĐỘNG VÀ SÓNG}
\subsubsection*{PHẦN I. Thí sinh trả lời câu hỏi từ câu 1 đến câu 18. Mỗi câu hỏi thí sinh chỉ chọn một phương án}
\setcounter{ex}{0}
	\begin{ex} %[3P1Y1-1][Loai1] 
		Một vật dao động điều hòa theo phương trình $x = A\cos\left(\omega t + \varphi \right)$ ($A > 0$). Biên độ dao động của vật là
		\choice
		{\True $A$}
		{$\varphi$}
		{$\omega$}
		{$x$}
		\loigiai{
			Biên độ dao động của vật là A.
		}
		%\haidong{5}
	\end{ex}
	\begin{ex} %[3P1Y1-2][Loai1] 
		Một chất điểm dao động có phương trình $x=10\cos\left(15t+\pi\right)$ (x tính bằng\ cm, t tính bằng s). Chất điểm này dao động với tần số góc là
		\choice
		{$20\ rad/s$}
		{$10\ rad/s$}
		{$5\ rad/s$}
		{\True $15\ rad/s$}
		\loigiai{
			Ta có: $\omega =15\ rad/s$.
		}	
		%\haidong{5}
	\end{ex}
	\begin{ex}%[3P1Y2-1][Loai1]
		Trong quá trình dao động, vận tốc của vật có giá trị bằng không (vật dừng lại tức thời) khi vật ở
		\choice
		{biên dương ($x = A$)}
		{biên âm ($x = -A$)}
		{đi qua vị trí cân bằng theo chiều âm}
		{\True biên dương hoặc biên âm}
		\loigiai{
			Trong quá trình dao động, vận tốc của vật có giá trị bằng không (vật dừng lại tức thời) khi vật ở biên dương hoặc biên âm.
		}
		%\haidong{5}
	\end{ex}
	%\loai{Dùng công thức độc lập x, v tìm A}
	\begin{ex}%[3P1K2-2][Loai1]
		Một chất điểm dao động điều hòa trên trục Ox. Trong thời gian $31,4\ s$ chất điểm thực hiện được $100$ dao động toàn phần. Lúc chất điểm đi qua vị trí có li độ $2\ cm$ theo chiều âm với tốc độ là $40\sqrt{3}\ cm/s$. Lấy $\pi = 3,14$. Biên độ dao động của chất điểm là
		\choice 
		{\True $4\ cm$}
		{$2\sqrt{2}\ cm$}
		{$3\ cm$}
		{$2\sqrt{3}\ cm$}
		\loigiai{
			\begin{itemize}
				\item Chu kì: $T = \dfrac{31,4}{100} = \dfrac{\pi}{10}\ s \Rightarrow \omega = 20\ rad/s$.
				\item Sử dụng phương trình độc lập thời gian, ta được phương trình:
				\begin{align*}
					\dfrac{2^2}{A^2} + \dfrac{(40\sqrt{3})^2}{20^2.A^2}= 1\Rightarrow A = 4\ cm.
				\end{align*}
			\end{itemize}
		}
		%\haidong{10}
	\end{ex}
	%\loai{Cho li độ. Tìm pha dao động}
	\begin{ex}%[3P1K4-1][Loai1]
		Cho một chất điểm đang dao động điều hòa với biên độ A. Biết tại thời điểm ban đầu $\left(t = 0\right)$ chất điểm đi qua vị trí có li độ $\dfrac{A}{2}$ và đang chuyển động theo chiều dương của trục tọa độ. Pha ban đầu của dao động là
		\choice 
		{$\dfrac{\pi}{4}\ rad$}
		{$\dfrac{\pi}{6}\ rad$}
		{$\dfrac{\pi}{3}\ rad$}
		{\True $-\dfrac{\pi}{3}\ rad$}
		\loigiai{
			Sử dụng đường tròn lượng giác ta suy ra được pha của dao động là $-\dfrac{\pi}{3}$.
		}
		%\haidong{10}
	\end{ex}
	%\loai{Khoảng thời gian ngắn nhất từ li độ $x_1$ đến li độ $x_2$ tính cả chiều chuyển động}
	\begin{ex}%[3P1B8-1][Loai1]
		Vật dao động điều hòa theo phương trình $x = 4\cos(8\pi t - \dfrac{\pi}{6})\ cm$. Thời gian ngắn nhất vật đi từ  $-2\sqrt{3}$  cm theo chiều dương đến vị trí có li độ  $2\sqrt{3}\ cm$ theo chiều dương là
		\choice
		{$\dfrac{1}{10}\ s$}
		{$\dfrac{1}{20}\ s$}
		{\True $\dfrac{1}{12}\ s$}
		{$\dfrac{1}{16}\ s$}
		\loigiai{
			\begin{center}
				\begin{tikzpicture}
					\tkzDefPoints{0/0/o, 5/0/a, 4.325/0/b, 3.525/0/c, 2.5/0/d,  -5/0/h, -4.325/0/g, -3.525/0/f, -2.5/0/e, -6/0/x', 6/0/x}
					\tkzDefShiftPoint[g](-90:0.5cm){g1}
					\tkzDefShiftPoint[o](-90:0.5cm){o1}
					\tkzDefShiftPoint[b](-90:0.5cm){b1}
					\tkzDrawSegments[blue,->,very thick](x',x)
					\tkzDrawPoints[size=5,fill=black](o,a,b,c,d,e,f,g,h)
					\tkzLabelPoint[above](o){\tiny $O$}
					\tkzLabelPoint[above](a){\tiny$A$}
					\tkzLabelPoint[above](b){\tiny$\dfrac{A\sqrt{3}}{2}$}
					\tkzLabelPoint[above](c){\tiny$\dfrac{A\sqrt{2}}{2}$}
					\tkzLabelPoint[above](d){\tiny$\dfrac{A}{2}$}
					\tkzLabelPoint[above](h){\tiny$-A$}
					\tkzLabelPoint[above](g){\tiny $-\dfrac{A\sqrt{3}}{2}$}
					\tkzLabelPoint[above](f){\tiny$-\dfrac{A\sqrt{2}}{2}$}
					\tkzLabelPoint[above](e){\tiny$-\dfrac{A}{2}$}
					\tkzLabelPoint[above](x){\tiny$x$}
					\tkzDrawSegments[->,very thick,red](g1,o1 o1,b1)
					\tkzLabelSegment[below](g1,o1){\tiny$\dfrac{T}{6}$}
					\tkzLabelSegment[below](o1,b1){\tiny$\dfrac{T}{6}$}
				\end{tikzpicture}
			\end{center}
			Khoảng thời gian ngắn nhất cần tìm là: $\Delta t  = \dfrac{T}{6}+\dfrac{T}{6}=\dfrac{T}{3}=\dfrac{1}{12}\ s$.
		}
		%\haidong{10}
	\end{ex}
	%\loai{Khoảng thời gian vật cách vị trí cân bằng lớn hơn b trong một chu kì}
	\begin{ex}%[3P1K8-3][Loai1]
		Cho một chất điểm đang dao động điều hòa với chu kì bằng $3\ s$ và biên độ bằng $6\ cm$. Trong một chu kì dao động, tổng thời gian mà khoảng cách từ chất điểm tới vị trí cân bằng lớn hơn $3\sqrt{2}\ cm$ là
		\choice
		{$\dfrac{4}{3}\ s$}
		{$\dfrac{2}{3}\ s$}
		{\True $1,5\ s$}
		{$\dfrac{3}{4}\ s$}
		\loigiai{
			\immini{\begin{itemize}
					\item 
					Ta có: $|x|> +3\sqrt{2}\ cm \Rightarrow x> +3\sqrt{2}\ cm$ hoặc $x< -3\sqrt{2}\ cm$
					\item Ta biểu diễn các điểm có li độ $3\sqrt{2}$ trên đường tròn lượng giác như hình vẽ.
					\item Trong một chu kì dao động, khoảng cách từ chất điểm tới vị trí cân bằng lớn hơn $3\sqrt{2}$  khi nó đi từ vị trí $M_1$ đến $M_2$ và từ $M_3$ đến $M_0\Rightarrow t = \dfrac{T}{4}+\dfrac{T}{4} = \dfrac{T}{2} = \dfrac{3}{2} = 1,5\ s$.
			\end{itemize}}{\begin{tikzpicture}[scale=1,thick,>=stealth']
					\tkzInit[ymin=-3,ymax=3,xmin=-3,xmax=3]\tkzClip
					\tkzDefPoints{-2.5/0/m, 2.5/0/n, 0/0/o, 0/2.5/p, 0/-2.5/q}
					\tkzDefShiftPoint[o](-135:2){0}
					\tkzDefShiftPoint[o](-45:2){1}
					\tkzDefShiftPoint[o](45:2){2}
					\tkzDefShiftPoint[o](135:2){3}
					\tkzDefPointBy[projection=onto m--n](1) \tkzGetPoint{h}
					\tkzDefPointBy[projection=onto m--n](0) \tkzGetPoint{h'}
					\draw(o)circle(2cm);
					\tkzDrawSegments[dashed](0,3 1,2)
					\tkzDrawSegments(p,q)
					\tkzDrawSegments[->](m,n)
					\tkzDrawSegments[blue,->](o,1 o,2 o,0 o,3)
					\tkzDrawPoints[size=3](o,1,h,0,2,3,h')
					\tkzMarkAngles[size=0.7cm,arrows=->](1,o,2)
					\tkzMarkAngles[size=0.7cm,arrows=->](3,o,0)
					\node at (0)[below left]{$M_0$};
					\node at (1)[below right]{$M_1$};
					\node at (2)[above right]{$M_2$};
					\node at (3)[above left]{$M_3$};
					\node at (2,0)[below right]{$6$};
					\node at (h')[below]{$-3\sqrt{2}$};
					\node at (h)[below]{$3\sqrt{2}$};
			\end{tikzpicture}}
		}
		%\haidong{15}
	\end{ex}
	%\loai{Thời gian vận tốc thỏa mãn khoảng cho trước trong một chu kì}
	\begin{ex}%[3P1K9-1][Loai1]
		Một vật dao động điều hòa với chu kì  là $\pi\ s$. Trong một chu kì, thời gian vận tốc của vật có giá trị không vượt quá một nửa tốc độ cực đại là
		\choice
		{$\dfrac{\pi}{6}$ s}
		{\True $\dfrac{2\pi}{3}$ s}
		{$\dfrac{\pi}{3}$ s}
		{$\dfrac{\pi}{4}$ s}
		\loigiai{
			\begin{itemize}
				\item	Dựa vào đường tròn vận tốc ta có thời gian thỏa mãn $v\le \dfrac{{{v}_{\max}}}{2}$ là $\dfrac{2T}{3}=\dfrac{2\pi}{3}$.
				
			\end{itemize}
		}
	\end{ex}
	%\loai{Tính quãng đường vật đi được trong quãng thời gian $\Delta t=nT$}
	\begin{ex}%[3P1KD-1][Loai1] 
		Một chất điểm dao động điều hoà với chu kì $T=\dfrac{\pi}{10}\ s$ và biên độ $A = 6\ cm$. Quãng đường vật đi được trong $10\pi\ s$ là
		\choice
		{$9\ m$}
		{\True $24\ m$}
		{$6\ m$}
		{$1\ m$}
		\loigiai{
			\begin{itemize}
				\item 		Ta có:
				$t = 100T \Rightarrow S = 100.4A = 400A = 24\ m$ (1T đi được 4A).
				
			\end{itemize}
			
		}
		%\haidong{15}
	\end{ex}
	\begin{ex}%[3P1KD-3][Loai1]
		Cho một chất điểm dao động điều hòa theo phương trình $x = A\cos(\pi t + \dfrac{\pi}{4})$, x tính bằng cm và t tính bằng s. Tính từ thời điểm ban đầu, tốc độ trung bình của vật trong $3,5\ s$ chuyển động là $6,36\ cm/s$. Biên độ dao động của chất điểm là
		\choice
		{$6\ cm$}
		{\True $3\ cm$}
		{$4\ cm$}
		{$5\ cm$}
		\loigiai{
			\immini{
				\begin{itemize}
					\item $T = \dfrac{2\pi}{\pi} = 2\ s$.
					\item Ta có: $t = 3,5\ s = 1.2 + 1,5$.
					\item Ta biểu diễn vị trí ban đâu của vật bằng điểm $M_0$ trên đường tròn lượng giác như hình vẽ.
					\item Trong khoảng thời gian 2 s, vật quay được góc $2\pi$, tương ứng đi được quãng đường 4A và lại quay về vị trí $M_0$.
					\item Trong khoảng thời gian 1,5 giây tiếp theo, vật quay thêm được 1 góc $\dfrac{3\pi}{2}$, đến vị trí $M_1$ trên đường tròn lượng giác.
					\item Quãng đường chất điểm đi được trong 3,5 s là $s = 4A \dfrac{A}{\sqrt{2}}+2A+ \dfrac{A}{\sqrt{2}} = 6A+A\sqrt{2} = 2,5.6,36\Rightarrow  A = 3\ cm$.
				\end{itemize}
			}{\begin{tikzpicture}[scale=1,thick,>=stealth']
					\tkzInit[ymin=-3,ymax=3,xmin=-3,xmax=3]\tkzClip
					\tkzDefPoints{-2.5/0/m, 2.5/0/n, 0/0/o, 0/2.5/p, 0/-2.5/q}
					\tkzDefShiftPoint[o](45:2){0}
					\tkzDefShiftPoint[o](-45:2){m'}
					\tkzDefPointBy[projection=onto m--n](m') \tkzGetPoint{h}
					\draw(o)circle(2cm);
					\tkzDrawSegments[dashed](m',0)
					\tkzDrawSegments(p,q)
					\tkzDrawSegments[->](m,n)
					\tkzDrawSegments[blue,->](o,0 o,m')
					\tkzDrawPoints[size=3](o,h,0,m')
					\tkzMarkAngles[size=0.7cm,arrows=->](0,o,m')
					\node at (0)[above right]{$M_0$};
					\node at (m')[below right]{$M_1$};
					\node at (2,0)[below right]{$A$};
					\node at (h)[below]{$\dfrac{A}{\sqrt{2}}$};
			\end{tikzpicture}}
		}
		%\haidong{15}
	\end{ex}
	%\loai{Cho đồ thị li độ tìm tần số, pha ban đầu, biên độ}
	\begin{ex}%[3P1-19.1K][Loai1] 
		\immini[thm]{
			Một vật dao động điều hòa trên trục Ox. Hình vẽ bên là đồ thị biểu diễn sự phụ thuộc của li độ x vào thời gian t. Tần số của dao động là	
			\choice
			{$\dfrac{5}{\pi}\ Hz$}
			{$2\ Hz$}
			{\True $2,5\ Hz$}
			{$\dfrac{2,5}{\pi}\ Hz$}
		}{
			\pgfplotsset{width=7.5cm,height=4cm,compat=1.4}
			\begin{tikzpicture}[draw=Mapcolor,scale=1,thick,>=stealth']
				\begin{axis}[
					axis lines=middle,
					xmin=0,xmax=4.5,xstep=1,ymin=-2.16,ymax=2.5,ystep=0.1,
					grid=major, %Có các tùy chọn: minor|major|both|none
					x grid style={dashed},
					y grid style={dashed},
					ytick={-2,0,2},
					xtick={1,2,3,4,5,6},
					xticklabels={,{0,2} },
					yticklabels={},
					xticklabel style={black,below},
					xlabel style={above=3pt},
					xlabel=$t\ (s)$,
					ylabel style={right=3pt},
					ylabel={$x\ (cm)$}
					]
					\addplot[line width=1.2pt,domain=0:4,samples=200,Mapcolor]{2*cos(deg(0.5*pi*x+pi/2))};
				\end{axis}
			\end{tikzpicture}
		}
		\loigiai{
			Chu kì của dao động: $T=0,4\ s\Rightarrow f=2,5\ Hz$.
		}
		%\haidong{15}
	\end{ex}	
	%\loai{Đồ thị pha dao động}
	\begin{ex}
		\immini[thm]{Hình bên là ảnh chụp một máy phát điện gió trên đất liền có chiều dài sải cánh là 120 m. Giả sử vào giữa trưa khi các tia sáng Mặt Trời vuông góc với mặt đất và cánh quạt đang quay ngược chiều kim đồng hồ với tốc độ ổn định là 20 vòng/phút. Chọn mốc tính thời gian lúc cánh quạt số 1 nằm song song với mặt đất, hướng theo chiều dương của trục Ox. Tốc độ bóng đen (theo đơn vị m/s) của đầu mút cánh số 2 trên mặt đất tại thời điểm $t=0,75\ s$ là
				\choice
		{$108\ m/s$}
		{$145\ m/s$}
		{\True $126\ m/s$}
		{$187\ m/s$}} 
		{\includegraphics[width=5.5cm]{Images/95528}}

		\loigiai{
			\immini{\begin{itemize}
					\item TLN: 126
					\item Ta có: $\omega=20\ \ct{vòng/phút}=\dfrac{2\pi}{3}\ rad/s$.
					\item $\Delta \varphi=\omega .\Delta t=0,75.\dfrac{2\pi}{3}=\dfrac{\pi}{2}$.
					\item Từ hình vẽ ta thấy li độ của đầu mút cánh số 2 ở thời điểm 0,75 s là: $x_2=-60\sqrt{3}\ cm$.
					\item Tốc độ khi đó: $v_2=\dfrac{\omega A}{2}=\dfrac{\dfrac{2\pi}{3}.120}{2}\approx 126\ m/s$.
			\end{itemize}}{\begin{tikzpicture}[scale=1,thick,>=stealth']
					\tkzDefPoints{0/0/o}
					\draw (o) circle(2cm);
					\draw[->] (-2.5,0) -- (2.5,0);
					\draw[very thick,red] (0,0) -- (2,0);
					\draw[very thick,red] (0,0) -- (-1,1.75);
					\draw[very thick,red] (0,0) -- (-1,-1.75);
					\draw[very thick,blue] (0,0) -- (0,2);
					\draw[very thick,blue] (0,0) -- (-1.75,-1);
					\draw[very thick,blue] (0,0) -- (1.75,-1);
					\tkzDrawPoints[size=4,fill=white](o)
					\node[red] at (2.25,2) {$(t=0)$};
					\node[red] at (2.5,0.5) {$(1)$};
					\node[red] at (-1.25,2) {$(2)$};
					\node[red] at (-1.25,-2.25) {$(3)$};
					\node[blue] at (0,2.25) {$(1)$};
					\node[blue] at (-2.25,-1.25) {$(2)$};
					\node[blue] at (2.25,-1.25) {$(3)$};
			\end{tikzpicture}}
		}	
		\haidong{15}
		
		
	\end{ex}
	%\loai{Từ tần số góc tính k hoặc m}
	\begin{ex} %[3P1B3-2][Loai1]  
		Một con lắc lò xo gồm lò xo có độ cứng $k=200\ N/m$, vật nhỏ khối lượng m, dao động điều hòa với tần số góc $20\ rad/s$. Giá trị của m là
		\choice
		{$100\ g$}
		{$200\ g$}
		{$400\ g$} 
		{\True $500\ g$}
		\loigiai{
			Tần số góc $\omega $ dao động của con lắc lò xo là:
			$\omega =\sqrt{\dfrac{k}{m}}\Leftrightarrow 20=\sqrt{\dfrac{200}{m}}\Rightarrow m=0,5\ kg=500\ g$.
		}	%<MyLT1>
		%\haidong{19}
	\end{ex}
	\begin{ex} %[3P1Y5-1][Loai1]  
		Một con lắc lò xo gồm vật nhỏ có khối lượng m và lò xo có độ cứng k, dao động điều hòa với phương trình $x=A\cos \left({\omega t+\varphi }\right)$. Mốc thế năng ở vị trí cân bằng. Cơ năng của con lắc là
		\choice
		{$\dfrac{1}{2}{mA}^2$}
		{\True $\dfrac{1}{2}{kA}^2$}
		{$\dfrac{1}{2}{mx}^2$}
		{$\dfrac{1}{2}{kx}^2$}
		\loigiai{
			Cơ năng của con lắc là $\dfrac{1}{2}{kA}^2$. 
		}
		%\haidong{15}
	\end{ex}
	\begin{ex} %[3P1YF-1][Loai1] 
		Một con lắc đơn có chiều dài $\ell$ dao động điều hòa tại nơi có gia tốc trọng trường $g$. Chu kì dao động riêng của con lắc này là
		\choice
		{$\dfrac{1}{2\pi}\sqrt{\dfrac{\ell}{g}}$}
		{$\dfrac{1}{2\pi}\sqrt{\dfrac{g}{\ell}}$}
		{\True $2\pi \sqrt{\dfrac{\ell}{g}}$}
		{$2\pi \sqrt{\dfrac{g}{\ell}}$}
		\loigiai{
			Chu kì dao động của con lắc đơn: $T=2\pi \sqrt{\dfrac{\ell}{g}}$}
		%\haidong{15}
	\end{ex}
	\begin{ex}%[3P1K6-2][Loai1] 
		Một con lắc lò xo dao động điều hòa theo phương thẳng đứng, trong quá trình dao động của vật lò xo có chiều dài biến thiên từ $12\ cm$ đến $20\ cm$. Biên độ dao động của vật là
		\choice
		{$8\ cm$}
		{\True $4\ cm$}
		{$16\ cm$}
		{$10\ cm$}
		\loigiai{
			Biên độ dao động của vật: $A=\dfrac{\ell_{\max}-\ell_{\min}}{2}=4\ cm$.
		}
		%\haidong{15}
	\end{ex}
	%\loai{Lí thuyết năng lượng}
	\begin{ex} %[3P1-16.1B][Loai1]   
		Với gốc thế năng được chọn tại vị trí cân bằng. Phát biểu nào sau đây là \textbf{sai} khi nói về cơ năng của con lắc đơn khi dao động điều hòa?
		\choice
		{Cơ năng bằng thế năng của vật ở vị trí biên}
		{Cơ năng bằng tổng động năng và thế năng của vật khi qua vị trí bất kì}
		{\True Cơ năng của con lắc đơn tỉ lệ thuận với biên độ góc}
		{Cơ năng bằng động năng của vật khi qua vị trí cân bằng}
		\loigiai{
			Cơ năng của con lắc đơn không tỉ lệ thuận với biên độ góc $\Rightarrow $ C sai.
		}
		%\haidong{15}
	\end{ex}
	\begin{ex}%[3P1KF-3][Loai1] 
		Con lắc đơn chiều dài $1\ m$ dao động nhỏ với chu kì $1,5\ s$ và biên độ góc là $0,05\ rad$. Độ lớn vận tốc khi vật có li độ góc $0,04\ rad$ là
		\choice
		{$9\pi \ cm/s$}
		{$3\pi \ cm/s$}
		{\True $4\pi \ cm/s$}
		{$\dfrac{4\pi}{3}\ cm/s$}
		\loigiai{Độ lớn vận tốc khi vật có li độ góc 0,04 rad
			\begin{align*}
				\begin{cases}
					T=2\pi \sqrt{\dfrac{\ell}{g}}\Rightarrow g=\dfrac{4\pi ^2\ell}{T^2}
					\\
					v^2=g\ell\left(\alpha_0^2-\alpha ^2\right)=\dfrac{4\pi ^2\ell^2}{T^2}\left(\alpha_0^2-\alpha ^2\right)\Rightarrow \left|v\right|=0,04\pi \ m/s\\ 
				\end{cases}
			\end{align*}
		}
		%\haidong{15}
	\end{ex}
\subsubsection*{PHẦN 2. Thí sinh trả lời từ câu 1 đến câu 4. Trong mỗi ý a), b), c), d) ở mỗi câu, thí sinh chọn đúng hoặc sai.}
\setcounter{ex}{0}
\begin{ex}
	Dao động cơ học xuất hiện trong nhiều hệ thống, từ dao động của con lắc đồng hồ đến sự nhún của lò xo giảm xóc ô tô.
	\choiceTF[t]
	{\True Dao động cơ học là chuyển động qua lại của một vật xung quanh vị trí cân bằng}
	{Dao động cơ của một vật luôn là dao động tuần hoàn}
	{\True Dao động điều hòa là dao động trong đó li độ được mô tả bằng một định luật dạng cosin (hay sin) theo thời gian}
	{\True Dao động điều hòa là trường hợp đặc biệt của dao động tuần hoàn}
	\loigiai{
		\begin{itemchoice}
			\itemch Đúng.
			\itemch Sai.
			\itemch Đúng.
			\itemch Đúng.
		\end{\itemchoice}
	}
	%\haidong{25}
\end{ex}
\begin{ex}
	Một chất điểm dao động điều hòa với phương trình $x=8 \cos 2 \pi t\ cm$.
	\choiceTF[t]
	{\True Chu kì dao động của chất điểm là $1 \ s$}
	{Trong nửa chu kì đầu tiên, vận tốc của vật bằng không ở thời điểm $t=\dfrac{1}{8} \ s$}
	{\True Thời điểm thứ nhất vật đi qua vị trí cân bằng là $\dfrac{1}{4} \ s$}
	{Từ thời điểm ban đầu, vật đi qua vị trí biên âm lần thứ $3$ là $4 \ s$}
	\loigiai{
		\begin{itemchoice}
			\itemch Đúng.
			\itemch Sai.
			\itemch Đúng.
			\itemch Sai.
		\end{\itemchoice}
	}
	%\haidong{24}
\end{ex}
\begin{ex}
	Một chất điểm dao động điều hòa dọc theo trục $O x$, quanh vị trí cân bằng $O$ với tần số $f$ và biên độ dao động là $A$.
	\choiceTF[t]
	{Thời gian ngắn nhất để chất điểm đi được quãng đường có độ dài $A \sqrt{3}$ là $\dfrac{\pi}{6 \omega}$}
	{Gọi $t_1$ và $t_2$ lần lượt là khoảng thời gian ngắn nhất và dài nhất để chất điểm đi được quãng đường bằng biên độ. Tỉ số $\dfrac{t_1}{t_2}$ bằng $\dfrac{\sqrt{2}}{2}$}
	{Thời gian ngắn nhất để chất điểm đi được quãng đường có độ dài $3A$ là $\dfrac{5}{6 f}$}
	{\True Thời gian ngắn nhất để chất điểm đi được quãng đường có độ dài $A$ là $\dfrac{1}{6 f}$}
	\loigiai{
		\begin{itemchoice}
			\itemch Sai.
			\itemch Sai.
			\itemch Sai.
			\itemch Đúng.
		\end{\itemchoice}
	}
	%\haidong{25}
\end{ex}
\begin{ex}
	Một con lắc lò xo dao động trên mặt phẳng nằm ngang với phương trình $x=5 \cos \left(5 \pi t-\dfrac{\pi}{4}\right)\ cm$. Khi vật ở vị trí cân bằng lò xo có chiều dài là $30 \ cm$.
	\choiceTF[t]
	{Biên độ dao động vật là $5 \pi\ cm$}
	{Tại vị trí $t=0$ vật có li độ $x=5 \sqrt{2}\ cm$}
	{\True Tốc độ cực đại trong quá trình dao động là $v=25 \pi\ cm/s$}
	{\True Chiều dài cực đại lò xo trong quá trình dao động là $35 \ cm$}
	\loigiai{
		\begin{itemchoice}
			\itemch Sai. Biên độ dao động vật là $5\ cm$.
			\itemch Sai. Tại vị trí $t=0$ vật có li độ $x= \dfrac{5\sqrt{2}}{2}\ cm$.
			\itemch Đúng.
			\itemch Đúng.
		\end{\itemchoice}
	}
	%\haidong{25}
\end{ex}
\subsubsection*{PHẦN 3. Thí sinh trả lời từ câu 1 đến câu 6.}
\setcounter{ex}{0}
\begin{ex}
	Một chất điểm dao động điều hòa dọc theo trục $O x$ trên quỹ đạo dài $10 \ cm$. Biết chất điểm thực hiện được 20 dao động toàn phần trong $5 \ s$. Tốc độ cực đại của vật trong quá trình dao động bằng bao nhiêu $m/s$ (làm tròn kết quả đến chữ số hàng phần trăm)?
	\shortans[oly]{1,26}
	\loigiai{
		$1,26 \ m/s$
	}
	%\haidong{32}
\end{ex}
\begin{ex}
	Một chất điểm dao động điều hòa theo phương trình $x=8 \cos \left(\dfrac{2 \pi}{3} t-\dfrac{\pi}{3}\right)$ ($x$ tính bằng $cm$, $t$ tính bằng $s$). Kể từ $t=11,25 \ s$, chất điểm cách vị trí cân bằng $4 \ cm$ và đang chuyển động ra xa vị trí cân bằng lần thứ 15 tại thời điểm bao nhiêu giây?
	\shortans[oly]{33} \loigiai{
		$33 \ s$
	}
	%\haidong{32}
\end{ex}
\begin{ex}
	Một vật dao động điều hòa với phương trình $x=4 \cos (\omega t-\dfrac{\pi}{3})\ cm$. Trong một chu kì dao động, khoảng thời gian mà vật có độ lớn gia tốc $\left|a\right|>\dfrac{a_{\max}}{2}$ là $0,4 \ s$. Khoảng thời gian kể từ khi vật dao động đến khi vật qua vị trí có tốc độ bằng $\dfrac{v_{\max}}{2}$ lần thứ hai là bao nhiêu giây?
	\shortans[oly]{0,15}
	\loigiai{
		0,15
	}
	%\haidong{30}
\end{ex}
\begin{ex}
	Một chất điểm dao động điều hòa với biên độ $4 \ cm$. Quãng đường lớn nhất mà chất điểm đi được trong $2 \ s$ là $12 \ cm$. Tần số góc của dao động bằng bao nhiêu rad/s (làm tròn kết quả đến chữ số hàng phần trăm)?
	\shortans[oly]{2,09}
	\loigiai{
		2,09
	}
	%\haidong{30}
\end{ex}
\begin{ex}
	Một chất điểm dao động điều hoà có hệ thức liên hệ giữa vận tốc $v$ và gia tốc $a$ là $\dfrac{v^2}{640}+\dfrac{a^2}{2,56}=1$ (trong đó, $v$ tính bằng $cm/s$ và $a$ tính bằng $m/s^2$). Biên độ dao động của chất điểm bằng bao nhiêu $cm$?
	\shortans[oly]{40}
	\loigiai{
		$40 \ cm$}
	%\haidong{30}
\end{ex}
\begin{ex}%[3P1K6-2][Loai1] 
	Con lắc lò xo treo thẳng đứng với đầu cố định của lò xo ở trên cao, quả nặng ở dưới thấp. Con lắc đang dao động điều hòa trên phương thẳng đứng, quanh vị trí cân bằng, với chu kì bằng $\dfrac{1}{\pi}\ s$. Biết độ nén lớn nhất của lò xo trong quá trình dao động là $2,5\ cm$. Cho gia tốc trọng trường $g = \pi^2 = 10\ m/s^2$. Tốc độ cực đại của vật trong lò xo trong quá trình dao động là bao nhiêu m/s?
	\shortans[oly]{1}
	\loigiai{
		\begin{itemize}
			\item Ta có: $\omega = \dfrac{2\pi}{T} = \dfrac{2\pi}{\dfrac{1}{\pi}} = 20\ rad/s$.
			\item Tại vị trí cân bằng, lò xo giãn một đoạn $\Delta \ell_0 = \dfrac{g}{\omega^2} = \dfrac{10}{20^2} = 0,025\ m = 2,5\ cm$.
			\item Mà độ nén lớn nhất của lò xo trong quá trình dao động là 2,5 cm $\Rightarrow A = 2,5 + 2,5 = 5\ cm\Rightarrow v_{\max} = \omega.A = 20.5 = 100\ cm/s = 1\ m/s$.
		\end{itemize}
	}
	%\haidong{30}
\end{ex}